\clearpage
\phantomsection
\section*{八、面向华为无线软件实现的技术方案选择与论证}
\addcontentsline{toc}{section}{八、面向华为无线软件实现的技术方案选择与论证}

\phantomsection
\subsection*{8.1 调研结论与工程前提}
\addcontentsline{toc}{subsection}{8.1 调研结论与工程前提}

面向华为无线软件实现，本章不应假设从零构建一套 LLM serving 系统。更现实的路线是基于已有Ascend推理生态进行策略增强：底层依托 Ascend NPU、CANN、torch-npu、ATB/MindIE Turbo 等软硬件栈；推理引擎优先考虑 vLLM Ascend 或 MindIE；上层再叠加业务感知 batching、状态感知 routing、移动性感知调度和验证环境中的策略控制。vLLM Ascend 官方文档明确将其定位为运行在 Ascend NPU 上的 vLLM 硬件插件，并说明这是 vLLM 社区支持 Ascend backend 的推荐方式；其安装文档也将 Linux、Ascend NPU、CANN、torch-npu、torch 和 NNAL 等列为基础依赖 \refcite{77}, \refcite{78}。因此，第八章的方案重点不应是“重新实现推理框架”，而应是“在现有推理框架可暴露的接口上做调度算法和系统集成”。

从能力边界看，vLLM Ascend 已经覆盖一批与本文算法方向直接相关的 serving 功能。其功能指南列出 dynamic batch、KV Cache Pool、KV Cache CPU offload、fine-grained tensor parallelism、layer sharding、speculative decoding、context parallel、dynamic chunked pipeline parallel 等能力入口 \refcite{79}；KV Cache Pool 文档进一步说明，在 PD disaggregation 场景下，Ascend Store 可管理 prefill 节点生成的 KV cache，并支持部分 decode 节点向 KV pool 写入状态 \refcite{80}；speculative decoding 文档列出 ngram、Medusa、EAGLE、MTP、draft model 等多种投机解码配置方式 \refcite{81}。这些能力说明，第 4、5 章中讨论的模型部署、batching、阶段解耦、KV 管理和 speculative verification，并非只能停留在论文原型层面，而是可以在Ascend推理栈中找到可实验的落点。

同时，MindIE 也是需要纳入工程路线的生产候选底座。Ascend社区的模型推理页面将 MindIE 定位为高性能 AI inference engine，强调其支持加速执行、调试、调优、迁移以及分层开放架构和统一接口；同一页面还提到开源推理引擎方向通过 vLLM 适配插件实现 NPU 与 vLLM 框架互通，MindIE Turbo 面向 NPU 上的 LLM 推理加速 \refcite{82}。因此，面向华为内部验证环境时，可以将 vLLM Ascend 作为开放生态和算法实验底座，将 MindIE/MindIE Service 作为更贴近Ascend生产栈的对照或集成目标。两者的共同点是：算法应尽量通过服务接口、调度参数、实例编排和观测数据接入，而不是直接改写底层算子或从头实现 runtime。

\phantomsection
\subsection*{8.2 基于 Ascend 推理栈的落地架构}
\addcontentsline{toc}{subsection}{8.2 基于 Ascend 推理栈的落地架构}

推荐架构应采用“现有推理引擎 + 轻量策略层 + 无线侧状态接入”的方式。推理实例层部署 vLLM Ascend 或 MindIE 服务，负责模型加载、连续批处理、KV 管理、token streaming、OpenAI-compatible API 或其他服务接口；策略层位于推理服务之外，负责读取请求 SLO、业务类型、用户入口、链路状态、实例负载和状态索引，并决定请求发送到哪个实例、是否启用特定模板、是否降级或回退；无线侧和边缘平台提供接入点、handover、链路质量、节点资源、容器状态和实例健康度等观测信息。

这种架构的核心原则是尽量少侵入推理引擎。单实例内的 token 调度、KV block 管理、算子执行和模型并行交给 vLLM Ascend 或 MindIE；外部策略层只做可解释、可回退的控制，例如设置请求优先级、选择目标实例、调整 admission 门限、选择部署模板、触发灰度切换或回退到云端。这样既能利用 Ascend 推理栈已有的性能优化，又能把本文综述的算法贡献落在无线软件可控制的层面。

对于验证环境，可以把系统拆成四个组件。第一是模型服务实例，运行 vLLM Ascend/MindIE，暴露 OpenAI-compatible、TGI-compatible 或内部 HTTP 接口，并上报 TTFT、TPOT、queue length、KV 水位、batch occupancy 等指标。第二是策略网关，接入无线侧请求，执行 SLO-aware routing、负载均衡、会话粘滞和异常回退。第三是状态索引服务，维护 session state、prefix cache 命中、KV 位置和模型副本信息；在早期阶段可以只维护元数据，不迁移真实 KV。第四是离线压测与画像模块，用 AISBench、vLLM benchmark 或内部压测工具生成模型画像、实例画像和链路画像，为在线策略提供参数边界。

\phantomsection
\subsection*{8.3 三类算法在现有框架中的落点}
\addcontentsline{toc}{subsection}{8.3 三类算法在现有框架中的落点}

业务感知 batching 的落点首先是推理实例内的调度参数和请求准入。vLLM 本身已经支持 OpenAI-compatible serving，并在官方文档中强调 PagedAttention、continuous batching、chunked prefill、prefix caching 等能力 \refcite{83}；vLLM Ascend 又提供 dynamic batch 和 Ascend 相关的 KV、并行与 speculative decoding 功能入口 \refcite{79}。因此，算法实现应优先围绕现有框架做三类增强：在请求入口按 SLO、prompt 长度和业务类型打标签；在实例侧记录 active batch、prefill/decode 队列和 KV 水位；在策略层根据观测结果调整 admission、等待窗口、实例权重和回退策略。若验收指标要求的 20\% rps 提升和 20\% E2E 时延下降无法仅通过外层策略达到，再有针对性地改造框架暴露的 scheduler 参数或队列接口。

请求调度与 task offloading 的落点更适合放在推理服务外部。vLLM Ascend 或 MindIE 实例可以被视为多个可选后端，策略网关根据用户入口、RTT、实例负载、模型能力、cache 命中和隐私约束选择目标实例。对于无线场景，不建议一开始就做跨实例 KV 迁移，而应先做 session 粘滞和状态元数据感知：如果用户会话已经在某个边缘实例上有上下文或 prefix cache，则优先保持该实例；如果用户 handover 后原实例链路变差，再比较远程访问、重新 prefill、转发到区域边缘或上云兜底的代价。这样可以先验证移动性信息对 routing 的收益，而不把工程风险集中到 KV 迁移接口上。

资源自适应模型部署的落点应从“动态任意切分模型”收敛为“Layer-wise 优先的拓扑感知混合切分部署”。在现有 Ascend 栈中，主线算法不应从 DP/TP/PP/EP 的笼统枚举开始，而应先围绕 Transformer layer、block 或 model shard 生成可部署的 pipeline stage：哪些连续层放在同一个 stage，stage 放到哪类 Ascend 节点，stage 之间经过哪条链路传输 activation 或 hidden state。只有在某个 stage 的显存或算力压力仍然过高时，再叠加 TP；如果目标模型包含 MoE，再考虑 EP；DP 则作为副本扩展和流量承载机制，而不是模型本体切分的核心。

因此，第 4 章的 taxonomy 在第 8 章中应被映射为一个分级实现路径。第一层是 layer-wise stage boundary、stage placement 和 pipeline balance，这是资源自适应模型切分算法的主要交付对象；第二层是 TP、context parallel、layer sharding 等框架可支持的补充并行能力，用于缓解单个 stage 内部的显存、算力或上下文压力；第三层是 EP、prefill/decode disaggregation、KV Cache Pool 和 speculative decoding 等扩展模板，用于 MoE、长上下文、阶段干扰或端边小模型可用的场景。系统先离线压测这些候选在不同模型、上下文长度和并发下的性能，再由在线策略在有限候选中选择可部署方案。需要深度改造的在线 operator split、跨实例真实 KV migration 或运行中频繁改变并行组，应作为风险项单独评估。

\phantomsection
\subsection*{8.4 风险、边界与建议}
\addcontentsline{toc}{subsection}{8.4 风险、边界与建议}

第一类风险是框架能力与论文假设之间的差距。论文中的动态模型切分、KV migration、state-aware routing 或 speculative verification 往往默认可以自由访问 runtime 内部状态；但在实际系统中，vLLM Ascend/MindIE 暴露的接口可能只支持一部分参数和指标。因此，报告后续应把算法分为“可直接交付”“需要轻量适配”“需要深度改造”三类。业务感知 admission、请求重排、外部 routing、实例权重调整、离线 layer profile、拓扑感知 layer-wise placement 和有限候选模板选择属于可直接交付或轻量适配；prefix 元数据索引、SLO-aware queue policy、chunked prefill 参数控制、TP/PP 模板画像和 PD/KV 扩展验证属于轻量适配；跨实例真实 KV 迁移、在线 operator split、运行中频繁重构并行组和深度 scheduler 改造则属于需要深度改造的风险项。

第二类风险是版本和硬件依赖。vLLM Ascend 文档显示，不同版本对 CANN、torch-npu、torch、NNAL 和 Atlas 设备有明确要求 \refcite{78}；同时，部分高级功能如 PD disaggregation、KV pool、speculative decoding、context parallel 或 layer sharding 可能与具体芯片、模型结构和版本绑定。实验章节必须记录软件版本、驱动版本、NPU 型号、模型版本和启动参数，否则结果很难复现，也不适合写成工程建议。

第三类风险是过早追求全局协同。模型切分、batching 和移动请求调度三项算法之间存在依赖关系，但项目交付并不要求一开始就实现全局最优控制面。更可控的做法是：T+3 到 T+8 月让模型切分和 batching 先形成稳定的可运行源码、画像数据和测试结果；T+9 到 T+12 月再把这些结果作为 routing 的实例能力输入，验证移动场景下的请求调度。这样既避免从零构建推理系统，又能让每一项算法都有清晰的基线、指标和回退方案。


